\begin{scanpage}{0037}{20}
where the integration contour passes above the pole $w=k\cos\phi_0$.

Let us note that integral (3.11) is a function of $|y|$, and let us change to polar coordinates according to the equations
\[\left.\begin{aligned}x&=r\cos\phi,\\|y|&=r\sin\phi\quad\text{with }\phi<\pi,\\|y|&=-r\sin\phi\quad\text{with }\phi>\pi.\end{aligned}\right\}\tag{3.12}\]
Furthermore, by carrying out the substitution
\[w=-k\cos\xi\quad(v=k\sin\xi),\tag{3.13}\]
we obtain
\[\left.\begin{aligned}I_1&=\frac1{i\pi k}\int_F\frac{e^{ikr\cos(\xi-\phi)}}{\cos\phi_0+\cos\xi}\,d\xi\quad\text{with }\phi<\pi,\\I_1&=\frac1{i\pi k}\int_F\frac{e^{ikr\cos(\xi+\phi)}}{\cos\phi_0+\cos\xi}\,d\xi\quad\text{with }\phi>\pi.\end{aligned}\right\}\tag{3.14}\]
The integration contour $F$ is shown in Figures 7a and 7b. In Figure 7a the cross-hatched areas indicate the sections in the plane of the complex variable $\xi$ in which $\operatorname{Im}\cos(\xi-\phi)>0$; in Figure 7b, the cross-hatched areas indicate the sections where $\operatorname{Im}\cos(\xi+\phi)>0$. Now let us deform the contour $F$ into the contour $G_1$ $(G_2)$ for the values $\phi<\pi$ $(\phi>\pi)$, and let us change to a new integration variable
\[\left.\begin{aligned}\zeta&=\xi-\phi&&\text{with }\phi<\pi,\\\zeta&=\xi-(2\pi-\phi)&&\text{with }\phi>\pi.\end{aligned}\right\}\tag{3.15}\]
As a result, we obtain the following expressions:
\[I_1=\frac1{i\pi k}\int_{D_0}\frac{e^{ikr\cos\zeta}\,d\zeta}{\cos\phi_0+\cos(\zeta+\phi)}+\begin{cases}0&\text{with }\phi>\pi-\phi_0,\\\dfrac2{k\sin\phi_0}e^{-ikr\cos(\phi+\phi_0)}&\text{with }\phi<\pi-\phi_0,\end{cases}\tag{3.16}\]
if $\phi<\pi$ and
\end{scanpage}
